\section{社会影响：白领通缩与失业窗口}

本期的另一半情绪来自社会影响。Coding 模型把顶尖研发人员生产力放大十倍到几十倍，研究和工程周期被压缩，组织自然会重新计算人力需求。白领通缩与失业窗口并不是远期科幻，而是生产力跃迁后组织吸收能力不足的现实问题。

\begin{figure}[H]
\centering
\includegraphics[width=0.92\textwidth]{social-impact-chain.png}
\caption{Coding 加速器如何传导到白领通缩与失业压力。}
\end{figure}

\begin{knowledgebox}{读图：社会冲击链条}
链条从 Coding 模型跃迁开始，传导到研发周期压缩，再到组织需求变化，最终形成工资和岗位压力。关键不是“AI 会不会替代人”这句空话，而是技术扩散速度是否超过组织和社会制度的吸收速度。
\end{knowledgebox}

\subsection{为什么是白领先感受到}

Coding、写作、分析、产品、运营、研究这些白领任务，很多都发生在数字世界，输入输出可被文本和代码表达，因此更容易被模型接入。蓝领和实体世界任务受机器人、硬件、环境复杂度限制，扩散速度可能更慢。

\begin{warningbox}{生产力提升不自动带来福利提升}
如果生产力收益集中在少数公司、少数顶尖个体或资本端，普通白领可能先经历工资压力和岗位减少。技术红利如何分配，是社会问题，不是模型能力自然会解决的问题。
\end{warningbox}

\subsection{投资新思考}

投资层面，本期暗含一个判断：若 Coding 是加速器，模型和工具链公司会重新估值；若模型成为 OS，应用入口和平台格局也会重排。但投资不能只看热词，而要看模型是否进入真实工作流、是否产生可持续反馈、是否有成本优势和分发入口。

\subsection{本章小结}

Coding/Agent 的社会影响来自生产力压缩。它会带来新公司、新入口和新投资机会，也会带来白领岗位重估、工资压力和制度挑战。


